% !TEX root = ../../tree_notes.tex
% =====================================================================
%  Chapter 12 -- Binary Lifting
%  Status: OUTLINE (not written yet). See structure/chapter.txt.
% =====================================================================
\bookchapter{Binary Lifting}{c12}
\chapterinfo{CF 1700--2100}{Outline}

\plannedtopic{Jumping by Powers of Two}{9}{%
  \item The $\text{up}[v][k] = 2^k$-th ancestor table
  \item Building the table in $O(n\log n)$
  \item Climbing to a given depth in $O(\log n)$
  \item $k$-th ancestor queries
}

\plannedtopic{Binary Lifting for LCA}{8}{%
  \item Equalise depths, then jump both nodes together
  \item The highest bit that keeps the nodes different
  \item LCA in $O(\log n)$ with a small constant
}

\plannedtopic{Lifting Arbitrary Path Data}{9}{%
  \item Store a value (min, max, sum, xor) per $2^k$ jump
  \item Combine on the way up to answer path queries
  \item Weighted edges and the distance to the LCA
  \item Jumps on a functional graph (successor lifting)
}

\plannedtopic{Binary Lifting Practice}{8}{%
  \item $k$-th ancestor of a node (LeetCode)
  \item Maximum edge on a path
  \item Binary lifting vs Euler+RMQ: when to use which
}

\closechapter
